\input{Datos.txt}
\documentclass[letterpaper,9pt]{article}

% *** Idioma ***
\usepackage[spanish,es-nodecimaldot]{babel}
\usepackage[utf8]{inputenc}

% *** Fuente ***
\usepackage{lmodern}
\usepackage[T1]{fontenc}

% *** Colores ***
\usepackage{xcolor}
\definecolor{dorado}{RGB}{190,150,91} 

% *** Margenes y Tamaño de Hoja
\usepackage{geometry}
\geometry{
left = 25mm,
right = 25mm,
top = 60mm,
bottom = 20mm,
headheight = 35mm
}

% *** Imagenes ***
\usepackage{graphicx}
\graphicspath{./imagenes}
\usepackage{float}
\usepackage{subcaption}
\usepackage{wrapfig}
\usepackage{tabularray}

% *** Fecha ***
\usepackage[useregional]{datetime2}


%No cortar palabras
\usepackage[none]{hyphenat}


% *** Encabezado ***
\usepackage{lastpage} 
\usepackage{fancyhdr}
\pagestyle{fancy}
\renewcommand{\headrulewidth}{0pt}
\lhead{
		\begin{tblr}{
		width = \linewidth,
		colspec = {Q[230]Q[507]Q[100]Q[100]},
		row{2} = {c},
		row{3} = {c},
		column{2} = {c},
		cell{1}{1} = {r=4}{c},
		cell{1}{2} = {r=3}{},
		cell{1}{3} = {c},
		cell{1}{4} = {c},
		cell{4}{3} = {c=2}{c},
                     cell{4}{2} = {c},
		vlines,
		hline{1,5} = {-}{},
		hline{2-3} = {3-4}{},
		hline{4} = {2-4}{},
	}
	\includegraphics[width=0.2\textwidth]{imagenes/Logo_de_la_Conagua} & {Comisión Nacional del Agua\\Gerencia de Tecnologías de Información \\y Comunicaciones\\ Subgerencia de Soporte Técnico, \\ Telecomunicaciones y Seguridad} & \tiny Hoja    & {\tiny \thepage\ de \pageref{LastPage}} \\
	&                                                                                       & \tiny Proyecto   & \tiny VPN \\
	&                                                                                       & \tiny Documento & {\tiny Formato }    \\
	& { Solicitud del Servicio de VPN}                                       &  \includegraphics[width=0.07\textwidth]{imagenes/Logo_Octubre_2024}       &           
\end{tblr}
}
\chead{

}
\rhead{}

% *** Pie de Pagina
%\renewcommand{\footrulewidth}{0.5pt}
\lfoot{}
\cfoot{}
\rfoot{}

\begin{document}
\sloppy



\begin{longtblr}[
	label = none,
	entry = none,
	]{
		width = \linewidth,
		colspec = {Q[203]Q[200]Q[200]Q[200]},
		column{2} = {c},
		cell{1}{3} = {r},
                     cell{2}{3} = {c},
		cell{2}{4} = {c},
		hlines,
		vlines,
	}
	Numero de reporte:  &          & Fecha:   & \today       \\
	Tipo de Movimiento: & Alta (\ALTA ) & Baja (\BAJA ) & Cambio (\CAMBIO ) 
\end{longtblr}

\vspace{-25pt}
\begin{longtblr}[
	label = none,
	entry = none,
	]{
		width = \linewidth,
		colspec = {Q[200]Q[300]Q[90]Q[90]},
		row{1} = {c},
                     row{1}={dorado},
                     cell{5}{3}={r},
                     cell{2}{2}={j},
                     cell{3}{2}={j},
		cell{1}{1} = {c=4}{0.874\linewidth},
		cell{2}{2} = {c=3}{0.649\linewidth},
		cell{3}{2} = {c=3}{0.649\linewidth},
		cell{4}{2} = {c=3}{0.649\linewidth},
		cell{6}{2} = {c=3}{0.249\linewidth},
		hlines,
		vlines,
	}
\textbf{Datos de Usuario (a)} &  &      &  \\
Nombre:                     &  \NOMBRE \\
Puesto o cargo: &   \PUESTO  \\
Unidad Administrativa:   &     \UA  \\
ID de Empleado:              & \ID  & Ext: &\EXT \\
Correo Institucional:         &  \CORREO  
\end{longtblr}

\vspace{-25pt}
\begin{longtblr}[
	label = none,
	entry = none,
	]{
		width = \linewidth,
		colspec = {Q[120]Q[160]Q[300]Q[100]},
		row{1} = {c},
                     row{1}={dorado},
                    cell{1}{1} = {c=4}{0.874\linewidth},		
		hlines,
		vlines,
	}
\textbf{Caracteristicas del equipo}         \\
\textbf{Marca:}       &  \MARCA & \textbf{\textbf{Cuenta con Anti-Malware:}}                           & Si (\SIANTI )~ No(\NOANTI ) \\
	\textbf{Modelo:}      &  \MODELO & {\textbf{Se encuentra vigente y actualizado:}\\\textbf{(Anti-Malware)}}    & Si (\SIVIGE )~ No(\NOVIGE ) \\
	\textbf{Serie:}       & \SERIE  & \textbf{Cuenta con S.O.:}                                  & Si (\SISO )~ No(\NOSO ) \\
	\textbf{Mac Address:} & \MACADDRESS  & {\textbf{Se encuentra licenciado y actualizado:}\\\textbf{(S.O.)}} & Si (\SILIC )~ No(\NOLIC ) 
\end{longtblr}

\vspace{-25pt}
\begin{longtblr}[
	label = none,
	entry = none,
	]{
                    	width = \linewidth,
		colspec = {Q[935]},
		cells = {c},
                     row{1}={dorado},
		hlines,
		vlines,
	}
	\textbf{Servicios que necesita acceder} \\\SERVICIOS  
\end{longtblr}

\vspace{-25pt}
\begin{longtblr}[
	label = none,
	entry = none,
	]{
		width = \linewidth,
		colspec = {Q[935]},
                     row{1}={dorado},
		cells = {c},
		row{2} = {},
		hlines,
		vlines,
	}
	\textbf{Justificación} \\\JUSTIFICACION
\end{longtblr}
\clearpage

La CONAGUA únicamente proveerá la habilitación del usuario dentro del concentrador VPN  proporcionando manual  para instalación y configuración del software necesario. Cualquier costo que se genere para hacer uso de este servicio, como la contratación de un servicio de Internet  o el uso de datos móviles, será de exclusiva responsabilidad del usuario.

\section*{Políticas de uso VPN}
{\small \begin{enumerate}
	\item Sólo los usuarios previamente autorizados podrán utilizar el servicio VPN y serán responsables del correcto uso del servicio de acceso remoto.
	\item Es responsabilidad del usuario asegurarse que ninguna otra persona utilice su cuenta de acceso, entendiendo que es de uso exclusivo para quienes se les ha asignado dichos privilegios.
	\item El uso del sistema VPN debe ser controlado utilizando una contraseña de autenticación fuerte  manteniéndola siempre en secreto. 
	\item El sistema VPN permitirá el tráfico de acuerdo con el perfil del usuario hacia y desde el equipo computacional a través del túnel VPN, el resto del tráfico pasará por la conexión respectiva.
	\item Solo se permite una sesión por usuario, por lo que no será posible conectarse simultáneamente en 2 o más dispositivos de computo.
	\item Todas los computadoras provistas por la institución conectados a las redes internas de la CONAGUA mediante VPN o cualquier otra tecnología deberán utilizar el software antivirus actualizado, provisto por el área de soporte del respectivo organismo. Para los equipos personales es responsabilidad del usuario proveer este software antivirus a sus equipos.
	\item La sesión será automáticamente desconectados de la sesión, una vez transcurrido 10 minutos de inactividad. El usuario deberá “loguearse” nuevamente para conectarse a la red CONAGUA. Procesos artificiales informáticos no deben ser utilizados para mantener la sesión abierta.
	\item Los usuarios de equipos computacionales que no son de la CONAGUA, deberán cumplir todas las disposiciones establecidas en las políticas  y además firmar un acuerdo de confidencialidad de la información.
	\item Al hacer uso de la tecnología VPN, los usuarios declaran conocer que sus computadoras, ya sea institucionales o personales son una extensión de las redes de la CONAGUA y como tales, están sujetos a las mismas normas y reglamentos que se aplican a los equipos dentro de esta institución.
\end{enumerate}}

\begin{table}[h]
	\centering
	\begin{tblr}{
			width = \linewidth,
			colspec = {Q[400]Q[400]},			
			row{1} = {c},
		}
		{Solicitó\\[1cm] \rule{6cm}{0.5mm}\\\NOMBRE\\\PUESTO} & {Autorizó\\[1cm] \rule{6cm}{0.5mm}\\\NOMBREJEFE\\\PUESTOJEFE} 
	\end{tblr}
\end{table}

\end{document}